\chapter*{Resumen}
\addcontentsline{toc}{chapter}{Resumen}
\markboth{Resumen}{Resumen}

La enseñanza del Método de Elementos Finitos (MEF) enfrenta una brecha pedagógica persistente: los textos clásicos presentan la formulación matricial con rigor pero de manera abstracta, mientras que el software profesional opera como una caja negra que oculta el procedimiento interno que el estudiante necesita comprender. Esta tesis aborda dicha brecha mediante el desarrollo de EduFEM, una herramienta educativa interactiva en español para el análisis estructural de problemas bidimensionales de elasticidad lineal en tensión plana y deformación plana, entendido como el análisis tenso-deformacional de medios continuos y no de sistemas estructurales discretos. El objetivo central es ofrecer un entorno donde el alumno construya el modelo, observe cada etapa del cálculo y verifique los resultados sin renunciar al rigor numérico. La metodología combina un motor de cálculo basado en elementos isoparamétricos cuadriláteros Q4 y Q9, una interfaz gráfica que integra pre-proceso, proceso y post-proceso sobre un único lienzo de malla, y un conjunto de módulos educativos que exponen paso a paso el mapeo isoparamétrico, el Jacobiano, las matrices de deformación y constitutiva, la rigidez elemental con cuadratura de Gauss, las fuerzas nodales equivalentes y el ensamblaje global. La corrección del motor se sometió a una batería de verificación y validación. Mediante el método de soluciones manufacturadas, en tensión plana y en deformación plana y sobre mallas regulares y distorsionadas, se confirmaron las tasas de convergencia asintóticas teóricas del desplazamiento —orden dos en norma $L^2$ y orden uno en seminorma $H^1$ para Q4, y orden tres y dos respectivamente para Q9— y la convergencia del campo de tensiones recuperado. La validación con la viga de Timoshenko arrojó, frente a la solución analítica, errores del 0,04\,\% en la tensión normal, del 0,26\,\% en la flecha y de hasta el 2,9\,\% en la tensión cortante; frente a un modelo de SAP2000, diferencias del 0,21\,\% en la tensión normal y de hasta el 0,56\,\% en desplazamientos. El caso de referencia de la membrana de Cook evidenció el bloqueo por cortante (\emph{shear-locking}) del elemento Q4 frente a la convergencia rápida de Q9, que para igual número de grados de libertad (GDL) reduce el error en torno a un orden de magnitud respecto del Q4. El aporte principal es EduFEM, una plataforma que articula teoría, implementación y experimentación numérica —incluida la generación automática de una memoria de cálculo paso a paso que el estudiante puede replicar y verificar a mano— para el aprendizaje del MEF.

\noindent\textbf{Palabras clave:} método de elementos finitos, software educativo, elementos isoparamétricos Q4/Q9, elasticidad plana, verificación y validación, memoria de cálculo